\subsection{גדלים אופייניים ותיקונים באטום המימן}
\begin{summarybox}
    {חזרה על גדלים יסודיים}בדיון על אטומים דמויי מימן (גרעין עם מטען $Ze$ ואלקטרון
    יחיד), הגדרנו שני גדלים סקלאריים מרכזיים:
    \begin{itemize}
        \item \textbf{רדיוס בוהר ($a_{0}$):} סקאלת האורך האופיינית.
            \begin{equation}
                a_{0} = \frac{4\pi\epsilon_{0} \hbar^{2}}{m_{e} e^{2}}\approx \SI
                {0.53}{\text{Å}}
            \end{equation}

        \item \textbf{אנרגיית רידברג ($Ry$):} סקאלת האנרגיה האופיינית.
            \begin{equation}
                Ry = \frac{m_{e} e^{4}}{32\pi^{2}\epsilon_{0}^{2} \hbar^{2}}\approx
                \SI{13.6}{eV}
            \end{equation}
    \end{itemize}ספקטרום האנרגיה נתון על ידי:
    \begin{equation}
        E_{n} = -\frac{Z^{2}}{n^{2}}Ry
    \end{equation}
\end{summarybox}
\subsubsection{תלות ב-$Z$ ורדיוס אופייני (תיקון טעות)}
בשיעור הקודם הוצג קשר שגוי לגבי הרדיוס הממוצע. הקשר הנכון הוא שככל שמטען הגרעין ($Z$)
גדל, המשיכה החשמלית מתחזקת והאטום "מתכווץ".הרדיוס הממוצע של האלקטרון (עבור המצב $n
=1$) הוא:
\begin{equation}
    \langle r \rangle \sim \frac{a_{0}}{Z}
\end{equation}האנרגיה הפוטנציאלית מתנהגת כמו
$\frac{Ze^{2}}{r}\sim \frac{Ze^{2}}{a_{0}/Z}\sim Z^{2}$, מה שמסביר את התלות הריבועית
ב-$Z$ באנרגיה.
\subsubsection{גבול הקירוב הלא-יחסותי}
המודל של שרדינגר הוא מודל לא-יחסותי. נשאלת השאלה: באילו תנאים קירוב זה נשבר?נבחן
דוגמה של אטום כבד מאוד, למשל אורניום ($Z=92$), עבור האלקטרון הפנימי ביותר ($n=1$).בדיקת
תקפות הקירוב הלא-יחסותי באורניום\textbf{נתונים:}
\begin{itemize}
    \item מספר אטומי: $Z \approx 100$ (לצורך הערכה סדר גודל).

    \item אנרגיית המנוחה של האלקטרון: $m_{e} c^{2} \approx \SI{511}{keV}$.
\end{itemize}\textbf{חישוב אנרגיית הקשר:} לפי הנוסחה הלא-יחסותית:
\begin{equation}
    |E_{1}| = Z^{2} \times \SI{13.6}{eV}\approx 100^{2} \times 13.6 = 136,000 \text{
    eV}= \SI{136}{keV}
\end{equation}

\textbf{מסקנה:} אנרגיית הקשר היא מסדר גודל של חמישית מאנרגיית המנוחה של האלקטרון
($\frac{136}{511}\approx 0.26$). משמעות הדבר היא שהאנרגיה הקינטית היא בסדר גודל
של אנרגיית המנוחה, ולכן המהירות האופיינית $v$ קרובה למהירות האור $c$. עבור
אטומים כבדים, השימוש במשוואת שרדינגר נאיבית אינו מדויק, ויש להשתמש בתיקונים
יחסותיים (משוואת דיראק).
\subsubsection{תיקון למסה מצומצמת (\LR{Reduced Mass})}
ההנחה שהגרעין נייח (מסה אינסופית) אינה מדויקת. עבור גרעין בעל מסה סופית
$M_{nuc}$, הבעיה הדו-גופית שקולה לבעיה של חלקיק אחד בעל \textbf{מסה מצומצמת}
$\mu$:
\begin{equation}
    \mu = \frac{m_{e} M_{nuc}}{m_{e} + M_{nuc}}
\end{equation}המעבר מהפתרון האידיאלי לפתרון הפיזיקלי מתבצע על ידי החלפת $m_{e}$
ב-$\mu$ בכל הנוסחאות:
\begin{itemize}
    \item $a_{0} \rightarrow a_{\mu}= a_{0} \frac{m_{e}}{\mu}$ (הרדיוס גדל מעט).

    \item $Ry \rightarrow Ry_{\mu}= Ry \frac{\mu}{m_{e}}$ (האנרגיה קטנה מעט).
\end{itemize}
\subsection{אלמנטי מטריצה וזוגיות (\LR{Parity})}
פונקציית הגל של אטום המימן נתונה על ידי:
\begin{equation}
    \psi_{nlm}(r, \theta, \phi) = R_{nl}(r) Y_{l}^{m}(\theta, \phi)
\end{equation}כאשר $Y_{l}^{m}$ הן ההרמוניות הספריות. נבחן אלמנטי מטריצה של
אופרטורים שונים.
\subsubsection{אופרטור סקלרי $r^{2}$}
האופרטור $r^{2}$ תלוי רק בקואורדינטה הרדיאלית ואינו משפיע על הזוויות. לכן, באלמנט
המטריצה תהיה הפרדה (פקטוריזציה):
\begin{align}
\mel{n l m}{r^2}{n' l' m'} &= \int_0^\infty r^2 dr
R_{nl}^*(r) r^2 R_{n'l'}(r) \times
\underbrace{\int d\Omega Y_{l}^{m*}(\Omega) Y_{l'}^{m'}(\Omega)}_{\delta_{ll'} \delta_{mm'}} \\
&= C(n,n',l) \delta_{ll'} \delta_{mm'}
\end{align}
האופרטור אינו משנה את המספרים הקוונטיים הזוויתיים.
\subsubsection{אופרטור וקטורי $z$ וכללי ברירה}
נבחן את אלמנט המטריצה של קואורדינטת הגובה $z = r\cos\theta$.
\begin{equation}
    \mel{n l m}{z}{n' l' m'}
\end{equation}משיקולי סימטריה לסיבוב סביב ציר $z$ (תלות ב-$\phi$), ומכיוון ש-$z$
אינו תלוי ב-$\phi$, נקבל מיד $\delta_{mm'}$.אולם, ניתן להסיק תובנה חזקה יותר משיקולי
זוגיות (\LR{Parity}).
\begin{definitionbox}
    {אופרטור הזוגיות $\mathcal{P}$}אופרטור הזוגיות מבצע שיקוף של הקואורדינטות
    דרך ראשית הצירים: $\vec{r}\rightarrow -\vec{r}$.תכונת ההרמוניות הספריות תחת זוגיות:
    \begin{equation}
        \mathcal{P}Y_{l}^{m}(\theta, \phi) = (-1)^{l} Y_{l}^{m}(\theta, \phi)
    \end{equation}הקואורדינטה $z$ היא אי-זוגית:
    $\mathcal{P}z \mathcal{P}^{\dagger} = -z$.
\end{definitionbox}
\begin{theorembox}
    {כלל ברירה לזוגיות}נבחן את התנהגות אלמנט המטריצה תחת זוגיות. נחדיר פעמיים את
    האופרטור ($\mathcal{P}^{\dagger} \mathcal{P}= \mathbb{I}$):
    \begin{align}
    \mel{nlm}{z}{n'l'm'}
    &= \mel{nlm}{\mathcal{P}^\dagger \mathcal{P} z \mathcal{P}^\dagger \mathcal{P}}{n'l'm'} \notag \\
    &= \mel{nlm}{\mathcal{P}^\dagger (-z) \mathcal{P}}{n'l'm'}
    \end{align}
    כאשר $\mathcal{P}$
    פועל על המצבים, הוא מוציא ערכים עצמיים של זוגיות:
    \begin{itemize}
        \item $\mathcal{P}\ket{n'l'm'}= (-1)^{l'}\ket{n'l'm'}$

        \item $\bra{nlm}\mathcal{P}^{\dagger} = \bra{nlm}(-1)^{l}$
    \end{itemize}
    נציב בחזרה:
    \begin{align}
    \mel{nlm}{z}{n'l'm'} &= (-1)^l (-1) (-1)^{l'} \mel{nlm}{z}{n'l'm'} \notag \\
    &= (-1)^{l+l'+1} \mel{nlm}{z}{n'l'm'}
    \end{align}
    כדי שאלמנט המטריצה לא יתאפס, נדרש ש:
    \begin{equation}
        (-1)^{l+l'+1}= 1 \implies l+l'+1 = \text{even}\implies l+l' = \text{odd}
    \end{equation}כלומר, המעבר מותר רק בין מצבים בעלי זוגיות הפוכה ($\Delta l$
    אי-זוגי). עבור דיפול חשמלי, כלל הברירה החזק יותר הוא $\Delta l = \pm 1$.
\end{theorembox}אותו טיעון תקף גם עבור $x$ ו-$y$, כיוון שגם הם משנים סימן תחת
זוגיות.
\subsection{סימטריות וחוקי שימור: ממשפט נתר הקלאסי לקוונטי}
\subsubsection{משפט נתר במכניקה קלאסית (חזרה)}
עבור מערכת המתוארת על ידי קואורדינטות $q_{a}$, משפט נתר קובע כי לכל סימטריה
רציפה של הלגרנז'יאן $L(q, \dot{q}, t)$ קיים גודל שמור.נניח טרנספורמציה
אינפיניטסימלית:
\begin{equation}
    q_{a} \rightarrow q_{a} + \epsilon f_{a}(q)
\end{equation}אם הלגרנז'יאן אינו משתנה ($\delta L = 0$), הגודל השמור הוא:
\begin{equation}
    Q = \sum_{a} \frac{\partial L}{\partial \dot{q}_{a}}f_{a} = \sum_{a} p_{a} f_{a}
\end{equation}
\subsubsection{קוונטיזציה של הגודל השמור}
במעבר למכניקה קוונטית, אנו הופכים את הגדלים לאופרטורים ($\hat{q}, \hat{p}$). כדי
לשמור על הרמיטיות של האופרטור $\hat{Q}$ (מכיוון ש-$\hat{p}$ ו-$\hat{f}$ לא בהכרח
חילופיים), נגדיר אותו באמצעות אנטי-קומוטטור (סימטריזציה):
\begin{equation}
    \hat{Q}= \frac{1}{2}\sum_{a} \left( \hat{p}_{a} \hat{f}_{a} + \hat{f}_{a} \hat
    {p}_{a} \right)
\end{equation}כעת נראה ש-$\hat{Q}$ הוא היוצר (\LR{Generator}) של הטרנספורמציה בתיאוריה
הקוונטית. נחשב את השינוי באופרטור הקואורדינטה $\hat{q}_{b}$ תחת הטרנספורמציה
האינפיניטסימלית המושרת על ידי $\hat{Q}$:
\begin{equation}
    \delta \hat{q}_{b} = \frac{i}{\hbar}\epsilon [\hat{Q}, \hat{q}_{b}]
\end{equation}
\begin{align}
[\hat{Q}, \hat{q}_b] &= \frac{1}{2} \sum_a [ \hat{p}_a \hat{f}_a +
\hat{f}_a \hat{p}_a , \hat{q}_b ] \notag \\
&= \frac{1}{2} \sum_a \left( [\hat{p}_a \hat{f}_a,
\hat{q}_b] + [\hat{f}_a \hat{p}_a, \hat{q}_b] \right)
\end{align}
נשתמש בזהות $[AB, C] = A[B,C
] + [A,C]B$. מכיוון ש-$\hat{f}_{a}$ היא פונקציה של הקואורדינטות בלבד, היא חילופית
עם $\hat{q}_{b}$ ($[\hat{f}_{a}, \hat{q}_{b}] = 0$). נותרנו רק עם הקומוטטורים של
התנע:
\begin{equation}
[\hat{p}_a, \hat{q}_b] = -i\hbar \delta_{ab}
\end{equation}
נציב זאת:
\begin{align}
[\hat{p}_a \hat{f}_a,
\hat{q}_b] &= \hat{p}_a [\hat{f}_a, \hat{q}_b] + [\hat{p}_a, \hat{q}_b] \hat{f}_a \notag \\
&= 0 + (-i\hbar \delta_{ab}) \hat{f}_a \\
[\hat{f}_{a} \hat{p}_{a}, \hat{q}_{b}] &= \hat{f}_{a} [\hat{p}_{a}, \hat{q}_{b}
    ] + [\hat{f}_{a}, \hat{q}_{b}] \hat{p}_{a} \notag \\
    &= \hat{f}_a (-i\hbar \delta_{ab}) + 0
\end{align}
בסכום הכולל:
\begin{equation}
[\hat{Q}, \hat{q}_b] = \frac{1}{2} \sum_a (-i\hbar \delta_{ab} \hat{f}_a
- i\hbar \delta_{ab} \hat{f}_a) = -i\hbar \hat{f}_b
\end{equation}
ולכן השינוי הוא:
\begin{equation}
    \delta \hat{q}_{b} = \frac{i\epsilon}{\hbar}(-i\hbar \hat{f}_{b}) = \epsilon
    \hat{f}_{b}
\end{equation}
\begin{resultbox}
    {מסקנה}האופרטור הקוונטי $\hat{Q}$, שנבנה מתוך הגודל השמור הקלאסי, מייצר
    בדיוק את הטרנספורמציה הצפויה על האופרטורים:
    \begin{equation}
        \hat{q}_{b} \rightarrow \hat{q}_{b} + \epsilon \hat{f}_{b}
    \end{equation}זהו קשר עמוק בין גדלים שמורים קלאסיים לבין יוצרים של סימטריות
    במכניקה הקוונטית (למשל: תנע הוא יוצר הזזה, המילטוניאן הוא יוצר קידום בזמן).
\end{resultbox}
\subsection*{בירורים נדרשים}
\begin{itemize}
    \item בשיעור נשאלה שאלה לגבי המקרה שבו $f_{a}$ תלוי ב-$\dot{q}$. המרצה ציין
        שהדבר אפשרי אך הדיון הוגבל לפונקציות של הקואורדינטות בלבד לצורך פשטות
        הקוונטיזציה (שכן $\dot{q}$ אינו מוגדר היטב כאופרטור בלתי תלוי).

    \item הייתה אי-בהירות מסוימת לגבי הסימון של הטרנספורמציה ההפוכה בחישוב הקומוטטורים,
        אך התוצאה הסופית $\epsilon \hat{f}_{b}$ עקבית עם התיאוריה.
\end{itemize}